\subsection{Raíz de Composición}

\subsubsection{Regla: los servicios se construyen solo en la raíz de composición}

\textbf{Regla:} Los servicios externos se crean con \texttt{new} o se registran en el contenedor únicamente en la raíz de composición: el punto de arranque de cada aplicación, como \texttt{Program}, las clases de registro de servicios del servidor o el método de inicialización del juego en el cliente. Ninguna otra clase crea los servicios que consume.

\textbf{Descripción:} La raíz de composición es el único lugar de la aplicación donde se arma el grafo de objetos; concentrarla ahí permite cambiar una implementación sin tocar las clases que la usan (van Deursen \& Seemann, 2019).

\textbf{Excepción:} Una clase puede crear las partes que forman parte de ella y comparten su tiempo de vida, como una unidad de trabajo que crea sus repositorios sobre el mismo contexto de base de datos.

\textbf{Código correcto:}
\begin{lstlisting}[style=csharpstyle]
public sealed class BoardGame : Game
{
    private TurnManager? _turnManager;

    protected override void Initialize()
    {
        _turnManager = new TurnManager(new DiceRoller());
        base.Initialize();
    }
}
\end{lstlisting}

\textbf{Código incorrecto:}
\begin{lstlisting}[style=csharpstyle]
public sealed class MatchScreen
{
    private readonly TurnManager _turnManager = new TurnManager(new DiceRoller());
}
\end{lstlisting}
